\section{Preliminaries}
\label{sec:prelim}


\begin{figure*}[t]
  \centering
    \begin{minipage}[c]{0.14\textwidth}
      \centering
      \includegraphics[width=\linewidth]{figures/03_loopmodel.pdf}
      \par\smallskip
      \textbf{(a)}
    \end{minipage}
    \hspace{0.02\textwidth}
    \begin{minipage}[c]{0.70\textwidth}
      \centering
      \includegraphics[width=\linewidth]{figures/01_looppipeline.pdf}
      \par\smallskip
      \textbf{(b)}
    \end{minipage} 
  \caption{
    \textbf{(a) Looped LLM architecture.}
    Blocks with dashed lines are optional.
    \textbf{(b) Comparing four decoding algorithms}
    (3 loops, shallow predictions after loop~1).
    Row $o_k$ is the position that predicts
    the $k$-th output token, with its input token on the left,
    and each column is one loop step.
    Autoregressive decoding starts a token after
    the previous one finishes.
    LoopSpec starts new token after one loop step, but the wrong draft
    (e.g. \texttt{the}) forces a restart.
    SuffixDecoding verifies multi-token drafts but waits
    for all three loops in every round.
    LoopFill (Section~\ref{sec:loopfill}) combines the two
    and takes 5 loop steps to finish, achieving $3\times$ speedup.
  }
  \label{fig:four-decoders}

  % Independent references without subcaption
  \begingroup
    \makeatletter
    \edef\@currentlabel{\thefigure(a)}
    \label{fig:loop-architecture}
    \edef\@currentlabel{\thefigure(b)}
    \label{fig:four-decoders-panel}
  \endgroup
\end{figure*}


\paragraph{Looped language models and pipelined decoding.}
A looped Transformer generates each token by repeatedly applying
a weight-shared block $f$ for $T$ iterations,
updating $h_i^{(t)}=f(h_{\le i}^{(t-1)})$
from an embedding $h_i^{(0)}$ of the previous token $o_{i-1}$.
Two properties enable efficient pipelined decoding.
First, intermediate hidden states can produce draft tokens
through the output head.
A draft token allows computation of the next token to begin
before the draft is verified.
Second, weight sharing allows computations for multiple
tokens at different loop depths to be processed together
in batches. These properties enable a decoding pipeline that
overlaps speculative drafting with full-depth verification.


We count one batched forward pass of the shared block $f$
as one loop step, during which each active token computation
advances by one loop. Consider LoopSpec~\cite{cho2026loopspecpipelinedselfspeculativedecoding}
with a draft depth $1\leq d<T$.
After $d$ loop steps, the model produces a draft token
$\hat o_i$ and uses it to start computing the next token
$o_{i+1}$, while continuing the computation of $o_i$
to depth $T$.
At depth $T$, the model verifies $\hat o_i$ and determines
the output token $o_i$.
As illustrated in Figure~\ref{fig:four-decoders-panel},
accepting the draft preserves the computation already
performed for $o_{i+1}$.
Otherwise, computations based on the rejected draft are
discarded, and computation of $o_{i+1}$ restarts using
the replacement token $o_i$.


\paragraph{Data-driven drafters.}
Model-free speculative decoding generates draft tokens from
previously text or decoding history without an auxiliary
draft model. LLMA and Prompt Lookup Decoding reuse continuations
from matching prompt spans~\cite{yang2023inferencereferencelosslessacceleration,saxena2023prompt}, ANPD and Token Recycling use $n$-grams and
cached token~\cite{ou-etal-2024-lossless,luo-etal-2025-turning}. 
SuffixDecoding~\cite{oliaro2025suffixdecodingextremespeculativedecoding}
extends this approach with suffix trees built from prompts and
historical outputs, adaptively constructing draft trees. The target LLM verifies the draft tree in parallel and
accepts the draft prefix with one
additional token by full-depth decoding.
In a looped model, each verification round takes $T$ loop steps,
and all token computations in the round start together.
The next round begins after the current round finishes.
In the SuffixDecoding panel of
Figure~\ref{fig:four-decoders-panel},
the three rounds commit 1, 3, and 1 tokens, respectively.
With $T=3$, generating these five tokens takes nine loop steps.


\paragraph{Breaking the pipeline ceiling.}
LoopSpec reduces the delay between consecutive token computations,
whereas SuffixDecoding allows multiple token computations to start
together. Vanilla LoopFill combines both mechanisms in a single draft tree,
adding suffix continuations before loop computation and shallow
continuations at depth $d$.

Let $N$ be the number of generated tokens, and let $\delta_i$ be
the number of loop steps between the start times of the computations
for $o_{i-1}$ and $o_i$ retained after verification.
The total loop-step cost and speedup over autoregressive decoding are
\begin{equation}
  K(N)=T+\sum_{i=2}^{N}\delta_i,
  \qquad
  S(N)=\frac{NT}{K(N)}.
  \label{eq:steps}
\end{equation}
Below, $S=\lim_{N\to\infty}S(N)$ denotes the average
loop-step speedup.

Let $\bar A\geq 1$ denote the average number of tokens accepted
per SuffixDecoding round, including the additional token produced
by full-depth verification. Let $\alpha\in[0,1]$ denote the
shallow-draft acceptance rate. For comparison, we assume that
LoopFill preserves the suffix coverage of SuffixDecoding and uses
the same $\alpha$ as LoopSpec.

In LoopSpec, each transition costs $d$ steps when its shallow
draft is accepted and $T$ steps otherwise, the average delay
is $\alpha d+(1-\alpha)T$.
In SuffixDecoding, transitions within a round have zero delay,
only a fraction $1/\bar A$ of transitions across drafting rounds cost $T$ steps.
LoopFill preserves the zero-delay transitions and replaces the
remaining $T$-step delays with $d$-step delays whenever a shallow
draft is accepted. Thus,
\begin{align}
  S_{\mathrm{LoopSpec}}
  &=\frac{T}{\alpha d+(1-\alpha)T}
    \leq\frac{T}{d},
  \label{eq:loopspec}\\
  S_{\mathrm{suffix}}
  &=\frac{T}{T/\bar A}=\bar A,
  \label{eq:suffix}\\
  S_{\mathrm{LoopFill}}
  &=\frac{T}{[\alpha d+(1-\alpha)T]/\bar A}
  \label{eq:combined}
\end{align}
The speedup from the two baseline methods are
\begin{equation}
  \begin{aligned}
    \frac{S_{\mathrm{LoopFill}}}{S_{\mathrm{suffix}}}
    &=\frac{T}{\alpha d+(1-\alpha)T}
      >1 \quad (\alpha>0),\\
    \frac{S_{\mathrm{LoopFill}}}{S_{\mathrm{LoopSpec}}}
    &=\bar A>1 \quad (\bar A>1).
  \end{aligned}
  \label{eq:relative-speedups}
\end{equation}
Relative to LoopSpec's upper nound speedup $T/d$, LoopFill achieves
\begin{equation}
  \begin{aligned}
    \frac{S_{\mathrm{LoopFill}}}{T/d}
    &=\frac{\bar A d}{\alpha d+(1-\alpha)T},\\
    S_{\mathrm{LoopFill}}>\frac{T}{d}
    &\quad\Longleftrightarrow\quad
      \bar A>\alpha+(1-\alpha)\frac{T}{d}.
  \end{aligned}
  \label{eq:pipeline-ceiling}
\end{equation}
In particular, when $\alpha=1$, LoopFill achieves
$S_{\mathrm{LoopFill}}=\bar A T/d$, exceeding the pipeline
ceiling whenever $\bar A>1$.